% GENERATED FILE. Edit canonical lessons, then run npm run generate:books.
% canonical-source-hash: fnv1a64:605e6b1ae8d73a8f
% canonical-lessons: LA-C151-superbus, LA-C151-benignus, LA-C151-crudelis, LA-C151-turpis, LA-C151-formosus

\chapter{Proud, Kind, Cruel, Ugly, Beautiful}
\label{ch:la-proud-kind-cruel-ugly-beautiful}
\hlchaptermodality{151}

\begin{chapteropening}
\textbf{By the end of this chapter:} \emph{I can say proud, kind, cruel, ugly and beautiful.}

Complete the last lesson of chapter 151: I can say proud, kind, cruel, ugly and beautiful.
\end{chapteropening}

\section[superbus, superba, superbum]{superbus, superba, superbum — proud}

\label{lesson:LA-C151-superbus}



\begin{quote}
Before the new one: say the Latin for timid, shy, then the Latin for bold.
\end{quote}

\subsection*{You'll want to know: superbus, superba, superbum}
\textbf{superbus, superba, superbum} — \textquotedblleft{}proud\textquotedblright{}.

\begin{cousinweb}
\textbf{superbus} is built on \emph{super}, \textquotedblleft{}above\textquotedblright{} (Lewis \& Short): one who sets himself above others. Greek \emph{hyperbios}, \textquotedblleft{}overbearing\textquotedblright{}, is a close parallel (Lewis \& Short).
\end{cousinweb}

\begin{grammarlens}[title={What you've built}]
One more word for this chapter.
\end{grammarlens}

\subsection*{Your turn}
\begin{itemize}
\raggedright
  \item \emph{Say it:} \emph{superbus}
  \item \emph{Say it:} \emph{superbus}, once more
  \item \emph{Say it:} say \emph{superbus}
  \item \emph{Recall:} say the Latin for timid, shy, then the Latin for bold, then say \emph{superbus} again
\end{itemize}

\begin{tcolorbox}[breakable,colback=teal!4,colframe=teal!35!black,title={Before you move on}]
What does superbus, superba, superbum mean? (\textquotedblleft{}Proud\textquotedblright{}.) Say it once more.
\end{tcolorbox}
\section[benignus, benigna, benignum]{benignus, benigna, benignum — kind}

\label{lesson:LA-C151-benignus}



\begin{quote}
Before the new one: say the Latin for bold, then the Latin for proud.
\end{quote}

\subsection*{You'll want to know: benignus, benigna, benignum}
\textbf{benignus, benigna, benignum} — \textquotedblleft{}kind\textquotedblright{}.

\begin{cousinweb}
\textbf{benignus} is usually explained as \emph{bene}, \textquotedblleft{}well\textquotedblright{}, plus the root of \emph{genus}, \textquotedblleft{}birth, kind\textquotedblright{} (Lewis \& Short), on the model of \emph{malignus}, \textquotedblleft{}ill-natured\textquotedblright{}: kindly by nature.
\end{cousinweb}

\begin{grammarlens}[title={What you've built}]
One more word for this chapter.
\end{grammarlens}

\subsection*{Your turn}
\begin{itemize}
\raggedright
  \item \emph{Say it:} \emph{benignus}
  \item \emph{Say it:} \emph{benignus}, once more
  \item \emph{Say it:} say \emph{benignus}
  \item \emph{Recall:} say the Latin for bold, then the Latin for proud, then say \emph{benignus} again
\end{itemize}

\begin{tcolorbox}[breakable,colback=teal!4,colframe=teal!35!black,title={Before you move on}]
What does benignus, benigna, benignum mean? (\textquotedblleft{}Kind\textquotedblright{}.) Say it once more.
\end{tcolorbox}
\section[crūdēlis, crūdēle]{crūdēlis, crūdēle — cruel}

\label{lesson:LA-C151-crudelis}



\begin{quote}
Before the new one: say the Latin for proud, then the Latin for kind.
\end{quote}

\subsection*{You'll want to know: crūdēlis, crūdēle}
\textbf{crūdēlis, crūdēle} — \textquotedblleft{}cruel\textquotedblright{}.

\begin{cousinweb}
\textbf{crūdēlis} is built on \emph{crūdus}, \textquotedblleft{}raw, bloody\textquotedblright{} (Lewis \& Short), usually traced to Proto-Indo-European \emph{*krewh2-}, \textquotedblleft{}raw flesh, blood\textquotedblright{} (EtymDB 2.1). English \emph{raw} is a cousin.
\end{cousinweb}

\begin{grammarlens}[title={What you've built}]
One more word for this chapter.
\end{grammarlens}

\subsection*{Your turn}
\begin{itemize}
\raggedright
  \item \emph{Say it:} \emph{crūdēlis}
  \item \emph{Say it:} \emph{crūdēlis}, once more
  \item \emph{Say it:} say \emph{crūdēlis}
  \item \emph{Recall:} say the Latin for proud, then the Latin for kind, then say \emph{crūdēlis} again
\end{itemize}

\begin{tcolorbox}[breakable,colback=teal!4,colframe=teal!35!black,title={Before you move on}]
What does crūdēlis, crūdēle mean? (\textquotedblleft{}Cruel\textquotedblright{}.) Say it once more.
\end{tcolorbox}
\section[turpis, turpe]{turpis, turpe — ugly, shameful}

\label{lesson:LA-C151-turpis}



\begin{quote}
Before the new one: say the Latin for kind, then the Latin for cruel.
\end{quote}

\subsection*{You'll want to know: turpis, turpe}
\textbf{turpis, turpe} — \textquotedblleft{}ugly, shameful\textquotedblright{}.

\begin{cousinweb}
\textbf{turpis} has no sure origin; neither Lewis \& Short's comparison nor the root that EtymDB 2.1 records is generally accepted.
\end{cousinweb}

\begin{grammarlens}[title={What you've built}]
One more word for this chapter.
\end{grammarlens}

\subsection*{Your turn}
\begin{itemize}
\raggedright
  \item \emph{Say it:} \emph{turpis}
  \item \emph{Say it:} \emph{turpis}, once more
  \item \emph{Say it:} say \emph{turpis}
  \item \emph{Recall:} say the Latin for kind, then the Latin for cruel, then say \emph{turpis} again
\end{itemize}

\begin{tcolorbox}[breakable,colback=teal!4,colframe=teal!35!black,title={Before you move on}]
What does turpis, turpe mean? (\textquotedblleft{}Ugly, shameful\textquotedblright{}.) Say it once more.
\end{tcolorbox}
\section[fōrmōsus, fōrmōsa, fōrmōsum]{fōrmōsus, fōrmōsa, fōrmōsum — beautiful}

\label{lesson:LA-C151-formosus}



\begin{quote}
Before the new one: say the Latin for cruel, then the Latin for ugly.
\end{quote}

\subsection*{You'll want to know: fōrmōsus, fōrmōsa, fōrmōsum}
\textbf{fōrmōsus, fōrmōsa, fōrmōsum} — \textquotedblleft{}beautiful\textquotedblright{}.

\begin{cousinweb}
\textbf{fōrmōsus} is built on \emph{fōrma}, \textquotedblleft{}shape, beauty\textquotedblright{} (Lewis \& Short), with \emph{-ōsus}, \textquotedblleft{}full of\textquotedblright{}: well-shaped. Where \emph{fōrma} comes from is uncertain; a link with Greek \emph{morphē}, \textquotedblleft{}shape\textquotedblright{}, is often suggested (EtymDB 2.1).
\end{cousinweb}

\begin{grammarlens}[title={What you've built}]
One more word for this chapter.
\end{grammarlens}

\subsection*{Your turn}
\begin{itemize}
\raggedright
  \item \emph{Say it:} \emph{fōrmōsus}
  \item \emph{Say it:} \emph{fōrmōsus}, once more
  \item \emph{Say it:} say \emph{fōrmōsus}
  \item \emph{Recall:} say the Latin for cruel, then the Latin for ugly, then say \emph{fōrmōsus} again
\end{itemize}

\begin{tcolorbox}[breakable,colback=teal!4,colframe=teal!35!black,title={Before you move on}]
What does fōrmōsus, fōrmōsa, fōrmōsum mean? (\textquotedblleft{}Beautiful\textquotedblright{}.) Say it once more.
\end{tcolorbox}
